\section{Design Rationale}
\label{sec:design-rationale}

The design rationale behind \SPHINCSA\ is to follow the original \SPHINCSP construction, apply the optimized \WOTSCS one-time signature scheme, replace the bottom few-time layer by BPORS+FP, and use a Diverse Layer Heights optimization for the hypertree. The idea behind \SPHINCSP was as follows. One can build a stateless hash-based signature scheme using a massive binary certification tree and selecting a leaf at random for each message to be signed. The problem with this approach is that the tree has to be extremely high, i.e., a height of about twice the security level would be necessary. This leads to impractical signature sizes. Using a hypertree instead of a binary certification tree allows the design to trade speed for signature size. However, this is still not sufficient by itself to get practical sizes and speed.

\paragraph{Mathematical hard problem.}
\SPHINCSA\ is a stateless hash-based post-quantum digital signature scheme. Its security is inherited from the hardness of generic attacks against the underlying hash and XOF primitives, following the \SPHINCSP proof framework. At the construction level, Kylin can be viewed as an improvement of the NIST-standardized SPHINCS+ design: it keeps the hash-based hypertree, replaces FORS by BPORS+FP, and replaces the Winternitz one-time signature layer with \WOTSCS.

\paragraph{Parameter selection criteria.}
The primary parameter-selection rule is to improve on comparable \SPHINCSP-style profiles while keeping the implementation and analysis close to the original construction. Under the fixed NGCC security levels, the search gives priority to smaller signatures, then faster verification, and then faster signing, while keeping the one-key signing budget at $q=2^{64}$. The selected instances in \cref{tab:ngcc-parameters} therefore use three main design levers: the constant-sum \WOTSCS encoding, the BPORS+FP parameters $(K,B,t,k,m_{\max})$, and the per-layer hypertree profile $(h_0,\ldots,h_{d-1})$.

Since one-time signatures are used extensively in the \SPHINCSP hypertree structure, optimizing the one-time signature size has a direct influence on the overall performance. With \WOTSCS, Kylin can retune the signature scheme to obtain smaller signatures while preserving comparable running time. BPORS+FP further reduces the bottom few-time payload by bucket-thinning child PORS keys and enforcing a fixed global authentication-node budget. Diverse Layer Heights removes the unnecessary restriction that $h$ be divisible by $d$: the layers may have different heights as long as their sum is $h$, which lets the search use balanced profiles such as $(7,7,7,7,7,7,7,8,8)$ for \texttt{Kylin-128}.

\paragraph{Failure probability and security impact.}
Kylin has no decryption- or decapsulation-style correctness failure event. For honestly generated keys and honestly generated signatures, verification is deterministic and accepts except for implementation faults or violation of the underlying cryptographic assumptions. The probability term used in the parameter search is therefore a forgery-success bound, not a correctness-failure probability.

For a signing budget of $q$ signatures under one key pair, the concrete
parameter-search exponent is derived from the BPORS+FP cover term in
\cref{eq:proof-generic-prob}.  The selected primary instances set $q=2^{64}$
and give the bit-security exponents shown in \cref{tab:ngcc-parameters}. Thus
the relevant probability term is bounded at the required NGCC levels for the
three primary instances. These exponents are used to judge the impact of
repeated bottom-key selection and generic forgery events on the claimed
security levels.

\paragraph{Other considerations.}
The repository also maintains $q=2^{24}$ companion profiles as reference-only additional implementations. These profiles are not the primary NGCC submission instances, but they are useful for applications with a smaller per-key signing budget and stricter signature-size preferences; see \cref{apx:q24-profiles}.

We are able to prove that the encoding scheme in \WOTSCS is size optimal among all encoding schemes over tree structures, even including those not necessarily efficiently computable. Moreover, to the best of our knowledge, no existing DAG-based scheme achieves better performance in terms of signature size, which renders our scheme the size-optimal one among all known schemes~\cite{optimalwots}. We are aware of probabilistic encoding schemes based on rejection sampling~\cite{cryptoeprint:2022/778}, and we consider those techniques to be of independent interest.

